\documentclass[10pt,a4paper]{amsart}
\usepackage{amsmath,amssymb,amsthm,graphicx}
\usepackage{verbatim}
\newenvironment{reference}{\comment}{\endcomment}
%\newenvironment{reference}{}{}
\newenvironment{slogan}{\comment}{\endcomment}
\newenvironment{history}{\comment}{\endcomment}
% For commutative diagrams we use Xy-pic
\usepackage[all]{xy}

% We use 2cell for 2-commutative diagrams.
\xyoption{2cell}
% Theorem environments.
%
\theoremstyle{plain}
\newtheorem{theorem}[subsection]{Theorem}
\newtheorem{proposition}[subsection]{Proposition}
\newtheorem{lemma}[subsection]{Lemma}

\theoremstyle{definition}
\newtheorem{definition}[subsection]{Definition}
\newtheorem{example}[subsection]{Example}
\newtheorem{exercise}[subsection]{Exercise}
\newtheorem{situation}[subsection]{Situation}

\theoremstyle{remark}
\newtheorem{remark}[subsection]{Remark}
\newtheorem{remarks}[subsection]{Remarks}

\numberwithin{equation}{subsection}

% Macros
%
\def\lim{\mathop{\mathrm{lim}}\nolimits}
\def\colim{\mathop{\mathrm{colim}}\nolimits}
\def\Spec{\mathop{\mathrm{Spec}}}
\def\Hom{\mathop{\mathrm{Hom}}\nolimits}
\def\Ext{\mathop{\mathrm{Ext}}\nolimits}
\def\SheafHom{\mathop{\mathcal{H}\!\mathit{om}}\nolimits}
\def\SheafExt{\mathop{\mathcal{E}\!\mathit{xt}}\nolimits}
\def\Sch{\mathit{Sch}}
\def\Mor{\mathop{\mathrm{Mor}}\nolimits}
\def\Ob{\mathop{\mathrm{Ob}}\nolimits}
\def\Sh{\mathop{\mathit{Sh}}\nolimits}
\def\NL{\mathop{N\!L}\nolimits}
\def\CH{\mathop{\mathrm{CH}}\nolimits}
\def\proetale{{pro\text{-}\acute{e}tale}}
\def\etale{{\acute{e}tale}}
\def\QCoh{\mathit{QCoh}}
\def\Ker{\mathop{\mathrm{Ker}}}
\def\Im{\mathop{\mathrm{Im}}}
\def\Coker{\mathop{\mathrm{Coker}}}
\def\Coim{\mathop{\mathrm{Coim}}}

% Boxtimes
%
\DeclareMathSymbol{\boxtimes}{\mathbin}{AMSa}{"02}

%
% Macros for moduli stacks/spaces
%
\def\QCohstack{\mathcal{QC}\!\mathit{oh}}
\def\Cohstack{\mathcal{C}\!\mathit{oh}}
\def\Spacesstack{\mathcal{S}\!\mathit{paces}}
\def\Quotfunctor{\mathrm{Quot}}
\def\Hilbfunctor{\mathrm{Hilb}}
\def\Curvesstack{\mathcal{C}\!\mathit{urves}}
\def\Polarizedstack{\mathcal{P}\!\mathit{olarized}}
\def\Complexesstack{\mathcal{C}\!\mathit{omplexes}}
% \Pic is the operator that assigns to X its picard group, usage \Pic(X)
% \Picardstack_{X/B} denotes the Picard stack of X over B
% \Picardfunctor_{X/B} denotes the Picard functor of X over B
\def\Pic{\mathop{\mathrm{Pic}}\nolimits}
\def\Picardstack{\mathcal{P}\!\mathit{ic}}
\def\Picardfunctor{\mathrm{Pic}}
\def\Deformationcategory{\mathcal{D}\!\mathit{ef}}
\usepackage{cleveref}
\crefname{originallemma}{original lemma}{original lemmas}
\Crefname{originallemma}{Original lemma}{Original lemmas}
\makeatletter
\newlabel{lemma-base-change-locally-quasi-finite}{{lemma-base-change-locally-quasi-finite}{source}{}{sourcepointer-1}{}}
\newlabel{lemma-base-change-locally-quasi-finite@cref}{{[originallemma][0][]lemma-base-change-locally-quasi-finite}{[1][0][]source}{}{}{}}
\makeatother
\crefname{originallemma}{original lemma}{original lemmas}
\Crefname{originallemma}{Original lemma}{Original lemmas}
\makeatletter
\newlabel{lemma-composition-finite-type}{{lemma-composition-finite-type}{source}{}{sourcepointer-2}{}}
\newlabel{lemma-composition-finite-type@cref}{{[originallemma][0][]lemma-composition-finite-type}{[1][0][]source}{}{}{}}
\makeatother
\crefname{originallemma}{original lemma}{original lemmas}
\Crefname{originallemma}{Original lemma}{Original lemmas}
\makeatletter
\newlabel{lemma-composition-separated}{{lemma-composition-separated}{source}{}{sourcepointer-3}{}}
\newlabel{lemma-composition-separated@cref}{{[originallemma][0][]lemma-composition-separated}{[1][0][]source}{}{}{}}
\makeatother
\crefname{originallemma}{original lemma}{original lemmas}
\Crefname{originallemma}{Original lemma}{Original lemmas}
\makeatletter
\newlabel{lemma-locally-quasi-finite-over-field}{{lemma-locally-quasi-finite-over-field}{source}{}{sourcepointer-4}{}}
\newlabel{lemma-locally-quasi-finite-over-field@cref}{{[originallemma][0][]lemma-locally-quasi-finite-over-field}{[1][0][]source}{}{}{}}
\makeatother
\crefname{originallemma}{original lemma}{original lemmas}
\Crefname{originallemma}{Original lemma}{Original lemmas}
\makeatletter
\newlabel{lemma-separated-implies-morphism-separated}{{lemma-separated-implies-morphism-separated}{source}{}{sourcepointer-5}{}}
\newlabel{lemma-separated-implies-morphism-separated@cref}{{[originallemma][0][]lemma-separated-implies-morphism-separated}{[1][0][]source}{}{}{}}
\makeatother
\crefname{originaltheorem}{original theorem}{original theorems}
\Crefname{originaltheorem}{Original theorem}{Original theorems}
\makeatletter
\newlabel{theorem-quasi-DM}{{theorem-quasi-DM}{source}{}{sourcepointer-6}{}}
\newlabel{theorem-quasi-DM@cref}{{[originaltheorem][0][]theorem-quasi-DM}{[1][0][]source}{}{}{}}
\makeatother
\crefname{originalreference}{original reference}{original references}
\Crefname{originalreference}{Original reference}{Original references}
\makeatletter
\newlabel{varieties-lemma-algebraic-scheme-dim-0}{{varieties-lemma-algebraic-scheme-dim-0}{source}{}{sourcepointer-7}{}}
\newlabel{varieties-lemma-algebraic-scheme-dim-0@cref}{{[originalreference][0][]varieties-lemma-algebraic-scheme-dim-0}{[1][0][]source}{}{}{}}
\makeatother
\title{A composite of locally quasi-finite morphisms of algebraic stacks is locally quasi-finite}
\author{The Stacks Project Authors}
\date{stacks-project:stacks-morphisms.tex:local-snapshot-20261012; GNU FDL 1.2 or later}
\begin{document}
\maketitle
\noindent\textit{Editorial note.} All delivered mathematical text is the author's own, copied from the pinned Stacks Project chapter source sources/stacks-algebraic-stacks-morphisms/pinned-original-source.tex; no statement or proof step was rewritten, reordered or removed.\par\medskip
\noindent\textit{Editorial note.} Preamble provenance: the Stacks Project keeps one single shared preamble for every chapter; this card takes it directly from the staged sources/stacks-algebraic-stacks-morphisms/preamble.tex as three verbatim spans (lines 8-12, 14-18 and 154-217), so no derived file was written and nothing was added or redefined. The wrapper class and the base packages are supplied by the assembler.\par\medskip
\noindent\textit{Editorial note.} Bibliography provenance: neither the counted statement nor its proof contains a citation (this lane rescans the whole span before writing the card), so the card carries no bibliography and no entry was reconstructed.\par\medskip
\noindent\textit{Editorial note.} Selection provenance: this card was opened under the file-based control contract (CONTROL.md), with the mechanical pre-card gate run on its statement and proof interval before the card was written.\par\medskip
\noindent\textit{Allowed standard prerequisites (editorial).}\begin{itemize}
\item The base change permanence of local quasi-finiteness. Invoked at the start of the counted proof (source lines 5045-5048). Makes the two morphisms after base change to a field locally quasi-finite.
\item The characterization of quasi-DM stacks through flat locally quasi-finite covers by schemes. Invoked inside the counted proof (source lines 5049-5056). Produces the scheme covering the middle stack.
\item The statement that a scheme locally quasi-finite over a field has discrete space. Invoked inside the counted proof (source lines 5057-5060). Gives discreteness of the presenting scheme.
\item The structure of a scheme locally quasi-finite over a field as a disjoint union of spectra of Artinian local algebras. Invoked near the end of the counted proof (source lines 5064-5074) as a cross-chapter pointer. Reduces discreteness of the fibre to the residue field points of the pieces.
\end{itemize}\medskip
\ifdefined\hypertarget\hypertarget{sourcepointer-1}{}\fi
\noindent\textit{Source reference.} Original lemma lemma-base-change-locally-quasi-finite: sources/stacks-algebraic-stacks-morphisms/pinned-original-source.tex, source lines 4989--4989. Navigation pointer only: the counted text refers to the Stacks item lemma-base-change-locally-quasi-finite, which is an existing labeled result of this same chapter and is not part of this card. The pointer carries the target's own label as its number and its definition line as its span; its content is not reproduced here.\par\medskip
\ifdefined\hypertarget\hypertarget{sourcepointer-2}{}\fi
\noindent\textit{Source reference.} Original lemma lemma-composition-finite-type: sources/stacks-algebraic-stacks-morphisms/pinned-original-source.tex, source lines 3488--3488. Navigation pointer only: the counted text refers to the Stacks item lemma-composition-finite-type, which is an existing labeled result of this same chapter and is not part of this card. The pointer carries the target's own label as its number and its definition line as its span; its content is not reproduced here.\par\medskip
\ifdefined\hypertarget\hypertarget{sourcepointer-3}{}\fi
\noindent\textit{Source reference.} Original lemma lemma-composition-separated: sources/stacks-algebraic-stacks-morphisms/pinned-original-source.tex, source lines 852--852. Navigation pointer only: the counted text refers to the Stacks item lemma-composition-separated, which is an existing labeled result of this same chapter and is not part of this card. The pointer carries the target's own label as its number and its definition line as its span; its content is not reproduced here.\par\medskip
\ifdefined\hypertarget\hypertarget{sourcepointer-4}{}\fi
\noindent\textit{Source reference.} Original lemma lemma-locally-quasi-finite-over-field: sources/stacks-algebraic-stacks-morphisms/pinned-original-source.tex, source lines 5002--5002. Navigation pointer only: the counted text refers to the Stacks item lemma-locally-quasi-finite-over-field, which is an existing labeled result of this same chapter and is not part of this card. The pointer carries the target's own label as its number and its definition line as its span; its content is not reproduced here.\par\medskip
\ifdefined\hypertarget\hypertarget{sourcepointer-5}{}\fi
\noindent\textit{Source reference.} Original lemma lemma-separated-implies-morphism-separated: sources/stacks-algebraic-stacks-morphisms/pinned-original-source.tex, source lines 1061--1061. Navigation pointer only: the counted text refers to the Stacks item lemma-separated-implies-morphism-separated, which is an existing labeled result of this same chapter and is not part of this card. The pointer carries the target's own label as its number and its definition line as its span; its content is not reproduced here.\par\medskip
\ifdefined\hypertarget\hypertarget{sourcepointer-6}{}\fi
\noindent\textit{Source reference.} Original theorem theorem-quasi-DM: sources/stacks-algebraic-stacks-morphisms/pinned-original-source.tex, source lines 4286--4286. Navigation pointer only: the counted text refers to the Stacks item theorem-quasi-DM, which is an existing labeled result of this same chapter and is not part of this card. The pointer carries the target's own label as its number and its definition line as its span; its content is not reproduced here.\par\medskip
\ifdefined\hypertarget\hypertarget{sourcepointer-7}{}\fi
\noindent\textit{Source reference.} Original reference varieties-lemma-algebraic-scheme-dim-0: sources/stacks-algebraic-stacks-morphisms/pinned-original-source.tex, source lines 4440--4440. Navigation pointer only: the counted text refers to the Stacks item varieties-lemma-algebraic-scheme-dim-0, which belongs to another chapter of the project and is not part of this card. The pointer carries the target's own label as its number because this chapter's staged snapshot does not contain the other chapter's tag table; its content is not reproduced here.\par\medskip
\begin{lemma}
\label{lemma-composition-locally-quasi-finite}
A composition of a locally quasi-finite morphisms is locally quasi-finite.
\end{lemma}

\begin{proof}
We have seen this for quasi-DM morphisms in
Lemma \ref{lemma-composition-separated}
and for locally finite type morphisms in
Lemma \ref{lemma-composition-finite-type}.
Let $\mathcal{X} \to \mathcal{Y}$ and $\mathcal{Y} \to \mathcal{Z}$
be locally quasi-finite. Let $k$ be a field and let
$\Spec(k) \to \mathcal{Z}$ be a morphism.
It suffices to show that $|\mathcal{X}_k|$ is discrete. By
Lemma \ref{lemma-base-change-locally-quasi-finite}
the morphisms $\mathcal{X}_k \to \mathcal{Y}_k$
and $\mathcal{Y}_k \to \Spec(k)$ are locally quasi-finite.
In particular we see that $\mathcal{Y}_k$ is
a quasi-DM algebraic stack, see
Lemma \ref{lemma-separated-implies-morphism-separated}.
By
Theorem \ref{theorem-quasi-DM}
we can find a scheme $V$ and a surjective, flat, locally finitely presented,
locally quasi-finite morphism $V \to \mathcal{Y}_k$. By
Lemma \ref{lemma-locally-quasi-finite-over-field}
we see that $V$ is locally quasi-finite over $k$, in particular
$|V|$ is discrete. The morphism
$V \times_{\mathcal{Y}_k} \mathcal{X}_k \to \mathcal{X}_k$ is
surjective, flat, and locally of finite presentation hence
$|V \times_{\mathcal{Y}_k} \mathcal{X}_k| \to |\mathcal{X}_k|$
is surjective and open. Thus it suffices to show that
$|V \times_{\mathcal{Y}_k} \mathcal{X}_k|$ is discrete.
Note that $V$ is a disjoint union of spectra of Artinian local
$k$-algebras $A_i$ with residue fields $k_i$, see
Varieties, Lemma \ref{varieties-lemma-algebraic-scheme-dim-0}.
Thus it suffices to show that each
$$
|\Spec(A_i) \times_{\mathcal{Y}_k} \mathcal{X}_k| =
|\Spec(k_i) \times_{\mathcal{Y}_k} \mathcal{X}_k| =
|\Spec(k_i) \times_\mathcal{Y} \mathcal{X}|
$$
is discrete, which follows from the assumption that
$\mathcal{X} \to \mathcal{Y}$ is locally quasi-finite.
\end{proof}

\end{document}
